
\usepackage{graphicx}%
\usepackage{multirow}%
\usepackage{amsmath,amssymb,amsfonts}%
\usepackage{amsthm}%
\usepackage{mathrsfs}%
\usepackage[title]{appendix}%
\usepackage{xcolor}%
\usepackage{textcomp}%
\usepackage{manyfoot}%
\usepackage{booktabs}%
\usepackage{algorithm}
\usepackage{algorithmicx}
\usepackage{algpseudocode}
\usepackage{multicol}
\usepackage{fancyvrb}
\usepackage{geometry}
\usepackage{titlesec}
\usepackage{xcolor}
\usepackage{parskip}
\usepackage[backend=biber,style=apa]{biblatex} % or use another style
\addbibresource{references.bib} % your .bib file
\usepackage{hyperref}
\usepackage{fontawesome5} % For the GitHub icon
\usepackage[justification=centering]{caption}
\usepackage[linesnumbered,ruled,vlined]{algorithm2e}
\setlength{\topmargin}{-0.4in}
\setlength{\textheight}{9in}
\usepackage{comment}
\usepackage{listings}%
% In your preamble:
\usepackage{enumitem} % For custom list formatting
\newlist{questions}{enumerate}{1}
\setlist[questions,1]{
  label=\textbf{Q\arabic*:}, % Label format "Q1:", "Q2:", ...
  leftmargin=2em,           % Indent spacing
  itemsep=1em               % Space between questions
}



%%%%


%% as per the requirement new theorem styles can be included as shown below
\theoremstyle{thmstyleone}%
\newtheorem{theorem}{Theorem}%  meant for continuous numbers
%%\newtheorem{theorem}{Theorem}[section]% meant for sectionwise numbers
%% optional argument [theorem] produces theorem numbering sequence instead of independent numbers for Proposition
\newtheorem{proposition}[theorem]{Proposition}%
%%\newtheorem{proposition}{Proposition}% to get separate numbers for theorem and proposition etc.

\theoremstyle{thmstyletwo}%
\newtheorem{example}{Example}%
\newtheorem{remark}{Remark}%

\theoremstyle{thmstylethree}%
\newtheorem{definition}{Definition}%

\raggedbottom

\usepackage{amssymb}

% Code Packages

\usepackage{listings}
\usepackage{xcolor}

\definecolor{headerbackground}{rgb}{1, 1, 1}
\definecolor{rulecolor}{rgb}{0.7, 0.7, 0.7}

% Define C keywords including standard library functions
\lstdefinestyle{algorithmstyle}{
  backgroundcolor=\color{white},
  basicstyle=\ttfamily\scriptsize,
  keywordstyle=\bfseries, % Bold for keywords
  commentstyle=\itshape, % Italic for comments
  stringstyle=,
  numberstyle=\tiny,
  numbers=left,
  numbersep=10pt,
  breaklines=true,
  frame=trbl,
  framerule=0.8pt,
  rulecolor=\color{rulecolor},
  framesep=5pt,
  xleftmargin=15pt,
  xrightmargin=5pt,
  language=C, % This is crucial
  tabsize=2,
  showstringspaces=false,
  alsoletter={\_},
  morekeywords={printf, scanf, malloc, free, sizeof, puts, gets}, % Add standard library functions
}

\newcommand{\codeheader}[1]{%
  \noindent%
  \fcolorbox{rulecolor}{headerbackground}{%
    \makebox[\dimexpr\linewidth-2\fboxsep-2\fboxrule][l]{%
      \textbf{Code:} \ttfamily{#1}%
    }%
  }%
  \vspace{-0.5\baselineskip}
}

\newcommand{\magn}[1]{
     \left[#1\right]
}